% English twin of the French section
\section{Consolidation, starter routines and transition to Year 12 (Première)}

\fichehead{Tue~1\textsuperscript{st}~June~2027}{Synthesis problems drawing on several topics}
\begin{fiche}
\qitem{Warm-up · Quick}{EQ}{Solve $3(x-1)=x+5$.}{$3x-3=x+5$, $2x=8$, $x=4$.}{}
\qitem{Core · Method}{GEO}{For which value of $k$ are the points $A(1;1)$, $B(3;5)$ and $C(k;9)$ collinear?}{$\vect{AB}\vc24$, $\vect{AC}\vc{k-1}{8}$: $2\times8-4(k-1)=0$, so $k=5$.}{}
\qitem{Core + · Method}{PRO}{A price increases by $10\,\%$ then decreases by $10\,\%$. What is the overall change?}{$1.1\times0.9=0.99$: a decrease of $1\,\%$.}{}
\qitem{Review · Quick}{PER}{Total surface area of a cuboid with dimensions $4$, $3$, $2$ (cm)?}{$2(4\times3+4\times2+3\times2)=52$ cm$^2$.}{}
\qitem{Thread · Quick}{SOL}{Calculate the exact volume of a sphere of radius $3$ cm.}{$V=\dfrac43\pi\times27=36\pi$ cm$^3$.}{}
\end{fiche}

\fichehead{Thu~3~June~2027}{Review of the six competences}
\begin{fiche}
\qitem{Warm-up · Metacognition}{MET}{List the six mathematical competences developed at lycée.}{Investigate, model, represent, reason, calculate, communicate.}{}
\qitem{Core · Metacognition}{MET}{Match each of the following to a competence: ``translating a problem statement into an equation''; ``sketching a curve''; ``justifying an equality with a written calculation''.}{Model; represent; reason (or communicate, depending on the emphasis placed on the write-up).}{}
\qitem{Core + · Metacognition}{MET}{To solve $2x+3=11$, which competence is mainly used?}{Calculate (manipulating a linear equation).}{For oral discussion: setting up an equation from a problem comes under ``model''.}
\qitem{Review · Problem}{MOD}{Water at $20$ °C is heated by $3$ °C per minute: $T(t)=20+3t$. After how many minutes does it reach $65$ °C?}{$20+3t=65\Leftrightarrow t=15$ min.}{}
\qitem{Thread · Quick}{EQ}{Solve in $\R$: $0.5x+1.2=2.7$.}{$x=3$}{}
\end{fiche}

\fichehead{Mon~7~June~2027}{Targeted reactivation of prerequisites for Première}
\begin{fiche}
\qitem{Warm-up · Quick}{DEV}{Expand $(x+3)^2$.}{$x^2+6x+9$.}{}
\qitem{Core · Method}{EQ}{Check that $x^2-5x+6=(x-2)(x-3)$ and deduce the solutions of $x^2-5x+6=0$.}{$(x-2)(x-3)=x^2-5x+6$; zero product: $x=2$ or $x=3$.}{}
\qitem{Core + · Anticipation}{MOD}{For $f(x)=x^2$, calculate the rate of change $\dfrac{f(3)-f(1)}{3-1}$.}{$\dfrac{9-1}{2}=4$: this prepares the idea of the derivative in Première.}{Première topic, simple consolidation calculation.}
\qitem{Review · Construction}{CON}{Place $M$ and $N$, the midpoints of $[AB]$ and $[AC]$, then draw $[MN]$. What can be said about $(MN)$ and $(BC)$? Calculate $MN$ if $BC=9$ cm.}{$(MN)\parallel(BC)$ and $MN=\dfrac92=4.5$ cm.}{}
\qitem{Thread · Write-up}{DEM}{Prove that a triangle with sides $7$, $24$ and $25$ is right-angled.}{$7^2+24^2=49+576=625=25^2$: converse of Pythagoras' theorem.}{}
\end{fiche}

\fichehead{Tue~8~June~2027}{Synthesis problems: choosing and justifying a strategy}
\begin{fiche}
\qitem{Warm-up · Reading code}{ALG}{What does $f(2)$ return?\par\smallskip{\ttfamily\small \hspace*{0.0em}def f(x):\\ \hspace*{2.4em}return x**2 - 3*x + 1\par}}{$4-6+1=-1$.}{}
\qitem{Core · Completion}{ALG}{Complete the code so that the function returns \code{True} when $A$, $B$, $C$ are collinear:\par\smallskip{\ttfamily\small \hspace*{0.0em}def alignes(xa, ya, xb, yb, xc, yc):\\ \hspace*{2.4em}return (xb-xa)*(yc-ya) - (yb-ya)*(xc-xa) == ....\par}}{We compare the determinant with $0$: \code{== 0}.}{}
\qitem{Core + · Completion}{ALG}{Complete the absolute value function:\par\smallskip{\ttfamily\small \hspace*{0.0em}def vabs(x):\\ \hspace*{2.4em}if x \textgreater{}= 0:\\ \hspace*{4.8em}return x\\ \hspace*{2.4em}else:\\ \hspace*{4.8em}return ....\par}}{\code{return -x}.}{Check: \code{vabs(-4)} returns $4$.}
\qitem{Review · Quick}{NUM}{Write $2.5\times10^3$ and then $6.2\times10^{-2}$ in decimal notation.}{$2\,500$; $0.062$.}{}
\qitem{Thread · Quick}{REP}{Give the reflection of $A(2;-3)$ in the origin.}{$(-2;3)$.}{}
\end{fiche}

\ifprof\else\clearpage\fi
\fichehead{Thu~10~June~2027}{Presentation of the specialist subjects; review}
\begin{fiche}
\qitem{Warm-up · Quick}{NUM}{Express $\dfrac38$ as a percentage.}{$37.5\,\%$.}{}
\qitem{Core · Metacognition}{MET}{What are your three revision priorities before Première?}{Personal answer: identify, in this year's sheets, the questions where you hesitated the most.}{}
\qitem{Core + · Quick}{EQ}{Solve $2x^2=18$.}{$x^2=9$: $x=3$ or $x=-3$.}{}
\qitem{Review · Trap}{PUI}{Calculate $(-2)^4$ and $-2^4$.}{$16$ and $-16$.}{The brackets decide what the exponent applies to.}
\qitem{Thread · True/false}{DIV}{True or false? Justify. If an integer is divisible by $4$ and by $6$, then it is divisible by $24$.}{False. Counter-example: $12$ is divisible by $4$ and by $6$, but not by $24$; it is divisible by $\ppcm(4\,;\,6)=12$.}{}
\end{fiche}
